\section{プログラミングの基礎}
\subsection{授業内容}
\subsubsection{科目の中でのこのコマの位置づけ}

本講義の主眼は「データから意味のある情報を取り出す」ことにある。すなわち，これまでのデータ駆動型(Data Driven)な発想を逆転し，データがどのように生成されたと考えるかという観点から，データ生成モデル(Data generating Model)による理解を試みる。

モデルは数学的表現がなされ，モデルの形成やデータとの照合(fitting)には計算機の利用が必須である。初回となる今回の目的は，プログラミングの基本的な考え方を身につけ，次回以降の本格的な運用に備えることである。
プログラミングの基礎はコマンドによる命令であり，コマンドの書き間違えはエラーとなって帰ってくる。一見不親切に思えるが，即時反応により学習の効率は良く，コード補完機能などを用いることで簡単なミススペルは回避することができる。小さなものから大きくしていくこと，1行ずつ実行することなど，基本的な姿勢についても理解する。

また，昨今ではchatGPTなど生成AIを用いたプログラミング・アシスト機能が発展しており，積極的な活用をすることで非常に効率よくプログラミングできる。
残念ながら，現状では「ハルシネーション(もっともらしい嘘)」の回答が与えられることも少なくない。あくまでも真偽の判断をするのは人間の側に残されており，参考にはできるものの正答が得られるとは限らないことに注意する。

\subsubsection{コマ主題細目}
\begin{description}
	\item[パソコンの基礎] スマートフォンやタブレットは「ファイル」を意識しないOSになっており，PCとそれらの境界はますますボーダレス化していくが，プログラミングを行う際はこれらの概念を知っておくことが重要である。PCの基本的な構造(BIOS/OS/アプリケーション)，ファイルの位置(パス，作業ディレクトリなど)などについての基礎知識を再確認する。
	
	\item[プログラミングの基礎] プログラミングにあたって重要なことは「思った通りに」動くのではなく，「書いた通りに」動くことである。ミススペルや大文字小文字の違いにも注意が必要である。コードのスペルチェックや補完機能を活用し，また1行ずつ確認しながら進めるといったプログラミングに関わる基本的な心構えについて理解する。

	\item[いくつかのプログラミング言語] Rはプログラミング言語の一種であると言っても良い。プログラミング言語には他にもPythonやBasic，C言語などがある。これらの基本的な関係について理解するとともに，コンパイラとインタプリタという実行形式の違いについても理解する。この違いは今後確率的プログラミング言語を利用する際の知識として生きてくる。

	\item[R言語の基本的な挙動] RおよびRStudioの最新版を導入し，改めて基本的な活用法を理解する。簡単な四則演算やオブジェクトへの代入，パッケージや関数の利用などを演習的に復習する。tidyverseという概念およびパイプ演算子をはじめとするプログラミング記法についても演習を通じて学ぶ。

	\item[生成AIの効率的な使い方] ChatGPTをコード生成時の補助役として活用する方法について理解する。ChatGPTは人間と違って，些細な質問を繰り返しても決して感情的になることはないため，質問の敷居は低い。一方で，間違えた答えを恥じることもなく，淡々と同じ間違いを重ねることもある。役割や方針を適切に与えた上で，問題を小分割してスモールステップでプログラミングを重ねていくことが基本である。また，結果が正しいかどうかは，必ず実際に動かしてユーザ側で判断する必要がある。

\end{description}
\subsubsection{キーワード}
\begin{itemize}
	\item プログラミング言語
	\item コンパイラ・インタプリタ
	\item オブジェクト指向プログラミング
	\item 生成AI
\end{itemize}
\subsection{授業情報}
\paragraph{コマの展開方法}
講義/遠隔/演習
\subsubsection{予習・復習課題}
\paragraph{予習}事前に環境の準備をしておく必要がある。環境の準備についてはいくつかの方策があり，これについては導入資料を参照しながら準備しておくこと。なお，環境準備中に問題が生じた場合はいち早く教員かTAに相談し，実行できるようにしておくこと。
\paragraph{復習}基本関数の使い方についての演習課題を課す。